%  LaTeX support: latex@mdpi.com 
%  For support, please attach all files needed for compiling as well as the log file, and specify your operating system, LaTeX version, and LaTeX editor.

%=================================================================
\documentclass[asi,article,submit,pdftex,moreauthors]{Definitions/mdpi} 

%=================================================================

% MDPI internal commands
\firstpage{1} 
\makeatletter 
\setcounter{page}{\@firstpage} 
\makeatother
\pubvolume{1}
\issuenum{1}
\articlenumber{0}
\pubyear{2023}
\copyrightyear{2023}
%\externaleditor{Academic Editor: Firstname Lastname}
\datereceived{10 May 2023} 
\daterevised{}
\dateaccepted{} 
\datepublished{} 
\hreflink{https://doi.org/} % If needed use \linebreak
%\doinum{}

%=================================================================
% Add packages and commands here. The following packages are loaded in our class file: fontenc, inputenc, calc, indentfirst, fancyhdr, graphicx, epstopdf, lastpage, ifthen, lineno, float, amsmath, setspace, enumitem, mathpazo, booktabs, titlesec, etoolbox, tabto, xcolor, soul, multirow, microtype, tikz, totcount, changepage, attrib, upgreek, cleveref, amsthm, hyphenat, natbib, hyperref, footmisc, url, geometry, newfloat, caption

\graphicspath{{figures/}} % path to folder with figures
\usepackage{subcaption}
\usepackage{listings}
\usepackage{xcolor}
\usepackage{gensymb} % for \textdegree
\usepackage{tabularx}
\usepackage{lipsum} 

%=================================================================
% Full title of the paper (Capitalized)
\Title{A GRASS GIS Framework for Detecting Changes in Patterns of the Ephemeral Salt Lakes Melrhir Chott and Merouane Chott, Algeria}

% MDPI internal command: Title for citation in the left column
\TitleCitation{A GRASS GIS Framework for Detecting Changes in Patterns of the Ephemeral Salt Lakes Melrhir Chott and Merouane Chott, Algeria}

% Author Orchid ID: enter ID or remove command
\newcommand{\orcidauthorA}{0000-0002-5759-1089} % Polina Add \orcidA{} behind the author's name
\newcommand{\orcidauthorB}{0000-0002-6461-1551} % Olivier Add \orcidB{} behind the author's name

% Authors, for the paper (add full first names)
\Author{Polina Lemenkova $^{1}$*\orcidA{}}

%\longauthorlist{yes}

% MDPI internal command: Authors, for metadata in PDF
\AuthorNames{Polina Lemenkova}

% MDPI internal command: Authors, for citation in the left column
\AuthorCitation{Lemenkova, P.}
% If this is a Chicago style journal: Lastname, Firstname, Firstname Lastname, and Firstname Lastname.

% Affiliations / Addresses (Add [1] after \address if there is only one affiliation.)
\address{%
$^{1}$ \quad Laboratory of Image Synthesis and Analysis (LISA), École polytechnique de Bruxelles (Brussels Faculty of Engineering), Université Libre de Bruxelles (ULB). Building L, Campus du Solbosch, ULB -- LISA CP165/57, Avenue Franklin D. Roosevelt 50, Brussels 1050, Belgium.
}

% Contact information of the corresponding author
\corres{Correspondence: polina.lemenkova@ulb.be; Tel.: +32 471 86 04 59 (P.L.)}

% Current address and/or shared authorship
%\firstnote{Current address: Affiliation 3.} 

% Abstract (Do not insert blank lines, i.e. \\) 
\abstract{
Automated classification of satellite images is a challenging task that enables to use remote sensing data for environmental modelling of the Earth's landscapes. In this document, we implement a GRASS GIS-based framework of land cover types discrimination to identify changes in the endorheic basins of the ephemeral salt lakes Chott Melrhir and Chott Merouane, Algeria, employing embedded algorithms of image processing. This study presents a dataset of the 10  Landsat 8-9 OLI/TIRS satellite images obtained from the USGS for a 10-year period of 2013-2023. The images were analysed for the changes of level in ephemerical lakes which experience temporal fluctuations in water level being dry for the most of the time and fed with water during rains. Two approaches are tested on the images for unsupervised classification for evaluating the algorithm of GRASS GIS: the 'i.cluster  module was used for generating image classes; the 'i.maxlik' was used for classification by the Maximal Likelihood approach. The document includes a technical description of the scripts used for image processing in the Appendix with comments on functionality in Methodology section. The results include the identified variations in the ephemeral salt lakes in the Algerian part of Sahara on 10-year period (2013-2023) using the series of the Landsat OLI/TIRS multispectral images by GRASS GIS. The main strength of the GRASS GIS framework is high speed, accuracy and effectiveness of the image processing for environmental monitoring. The presented GitHub repository containing GRASS GIS scripts used in this study provides a reference to satellite image interpretation which can be used in similar studies on environmental monitoring of the arid and semi-arid areas of Africa.
}

% Keywords
\keyword{Sahara Desert; Africa; salt lake; cartography; image analysis; satellite image; remote sensing; geography; GRASS GIS; mapping} 

% The fields PACS, MSC, and JEL may be left empty or commented out if not applicable
\PACS{91.10.Da; 91.10.Jf; 91.10.Sp; 91.10.Xa; 96.25.Vt; 91.10.Fc; 95.40.+s; 95.75.Qr; 95.75.Rs; 42.68.Wt}
\MSC{86A30; 86-08; 86A99; 86A04}
\JEL{ Y91; Q20; Q24; Q23; Q3; Q01; R11; O44; O13; Q5; Q51; Q55; N57; C6; C61}

%%%%%%%%%%%%%%%%%%%%%%%%%%%%%%%%%%%%%%%%%%
\begin{document}

\section{Introduction}

\subsection{Background and Motivation}
Detecting changes in water level of the ephemeral salt lakes using remote sensing data is an important task in hydrological analysis of arid regions, as it is often the basis of higher level applications, such as environmental monitoring, water resource management, climate change studies, analysis of groundwater distribution, etc. Located in the northern Africa's Maghreb region in Sahara desert, salt lakes, or "chotts", are subject to temporal fluctuations in water level throughout the calendar year. Chotts remain dry for the most of the time but are fed by water from wadi during winter or occasional rains. Monitoring such changes in water supply using variability of contours of the ephemeral stream beds is possible using satellite image classification which is based on different spectral reflectance of the objects visible on the land surface \cite{Medjani}. Changes in lake contours are used as indicators for monitoring land cover changes and hydrological modelling. 

Satellite images are a key data source in computational ecology and environmental monitoring. Various colours and brightness of the land cover classes visible on the images are determined by the spectral reflectance of the pixels and can be used for mapping the landscapes of the Earth \cite{Takorabt,Laghouag,Lemenkova202227,info14040249,Amri}. Using remote sensing data enables to detect heterogeneity in  the Earth's landscapes \cite{Abbas}, analyse geological processes  \cite{ARAIBIA2022104455,LAMRI2016143}. Furthermore, remote sensing data are used to record geologic processes \cite{Deroin2019,data7060074}, or support environmental analysis \cite{Lemenkova202229}. Further, spectral reflectance of the land cover types identified on the satellite image may be used to detect soil salinity \cite{abdennour2019detection}. Finally, remote sensing data are applicable in geomorphological studies and can be used to reflect terrain types \cite{Rebillard,technologies11020046,Lemenkova202301}, and climate studies such as mapping evapotranspiration and droughts \cite{ijgi11090473}, and characteristics of soil layer \cite{MULDER20111}.

This makes satellite images as a data source in many applications of Earth science such as mapping distribution of dust in arid regions \cite{Parajuli}, assessing and mapping erosion \cite{Mihi}, aquifer characterisation \cite{AOUDIA2020103920}, evaluating groundwater potential \cite{Derdour}, geological investigations of rock physics, tectonics and hydrogeological simulations \cite{NEDJRAOUI2021104348}. The analysis of the satellite images can also be used for assessing of the geophysical setting in the desert areas such as African Sahara by integrated geological and geospatial data modelling and analysis \cite{ZERROUK2017146}. An important consideration for the use of satellite images for cartographic data analysis is that they have to be classified into maps of land cover types for interpretation and analysis of landscapes. However, manual classification using traditional GIS is a time-consuming process involving many iterations which might lead to errors and results in low accuracy. The problem is therefore focused on finding the effective algorithms to classification of satellite images.

Landsat 8-9 OLI/TIRS images recorded since 2013 as a continuation of the Landsat data collection form a valuable pool of big Earth geospatial data to evaluate land cover types and landscapes of the Earth. However, such datasets should be processed and classified to enable the use of the satellite images in the environmental applications. The use of traditional GIS software is a time‐consuming and cumbersome work which can lead to misclassification of pixels during identification of land cover types. To address these limitations, we suggest using GRASS GIS scripts to automate classification of the remote sensing data by employing algorithms of image processing and geospatial analysis for discrimination of land cover types and detecting changes in the north-east Algeria. Besides the existing applications in the geomorphometry \cite{HOFIERKA2009387}, it can also be used for satellite image processing for environmental applications \cite{jmse11040871}.

The need to automate classification of the satellite images in order to build a structured workflow for remote sensing data processing intensifies the need to develop tools that allow the automated image processing with reduced workload. Recently, an existing types of classification methods (clustering, maximal likelihood, K-means) have revealed the importance of the optimisation of the algorithms of image processing. Scripting techniques of the GRASS GIS ensure improvement of the accuracy and accelerates the speed of classification. 

\subsection{Study Area}

The environmental analysis of the Algerian chotts have resulted a variety of the existing papers, e.g., geochemical analysis of the soil salinity \cite{Abdennour}, hydrogeological studies to evaluate geothermal resources \cite{BENMARCE2021104285,MELOUAH2021102211}, analysis of the mineralogy of desert alluvial soils in wadi of northern Sahara \cite{Djili}, brine geochemistry and mineralogy of the Algerian playa \cite{HamdiAissa,Valles}, or hydrological reconstruction in chotts \cite{Samie}. 

\begin{figure}[H]
\begin{adjustwidth}{-\extralength}{0cm}
%\centering
	\hspace{60pt}\includegraphics[width=16.0 cm]{Fig_01.jpg}
\end{adjustwidth}
\caption{Topographic map of Algeria. Software: GMT v. 6.1.1. Data sources: GEBCO and DCW. Rotated red square indicates study area with Chotts Melrhir and Merouane. Map source: authors.
\label{fig01}}
\end{figure}

%%%%%%%%%%%%%%%%%%%%%%%%%%%%%%%%%%%%%%%%%%
\section{Materials and Methods}

%----------- subsection ----------->
\pagebreak
\subsection{Data}

Figure \ref{fig:03} shows the dynamics visible on the Landsat satellite scenes for the 10 consecutive years 2014-2023 in January on the following dates: \textbf{a}) 2014/01/01, (\textbf{b}) 2015/01/04, (\textbf{c}) 2016/01/07, (\textbf{d}) 2017/01/25, (\textbf{e}) 2018/01/28, (\textbf{f}) 2019/01/15, (\textbf{g}) 2020/01/02, (\textbf{h}) 2021/01/04, (\textbf{i}) 2022/01/15, (\textbf{j}) 2023/01/10, (\textbf{k}) Enlarged fragment of the Melrhir Chott and Merouane Chott on the Google Maps (2023), (\textbf{l}) Enlarged fragment of the Melrhir Chott and Merouane Chott on the Google Earth aerial scene (2023).

\begin{figure}[H]
\makebox[0.90\linewidth][r]{%
	\begin{subfigure}[b]{.40\textwidth}
		\centering
			\includegraphics[width=0.87\textwidth]{Fig_03a.jpg}
		\caption{2014}
	\end{subfigure}%
	\begin{subfigure}[b]{.40\textwidth}
		\centering
		\includegraphics[width=0.87\textwidth]{Fig_03b.jpg}
		\caption{2015}
	\end{subfigure}%
	\begin{subfigure}[b]{.40\textwidth}
		\centering
		\includegraphics[width=0.87\textwidth]{Fig_03c.jpg}
		\caption{2016}
	\end{subfigure}%
}\\
\vfill \vspace{1mm}
\makebox[0.90\linewidth][r]{%
	\begin{subfigure}[b]{.40\textwidth}
		\centering
			\includegraphics[width=0.87\textwidth]{Fig_03d.jpg}
		\caption{2017}
	\end{subfigure}%
	\begin{subfigure}[b]{.40\textwidth}
		\centering
		\includegraphics[width=0.87\textwidth]{Fig_03e.jpg}
		\caption{2018}
	\end{subfigure}%
	\begin{subfigure}[b]{.40\textwidth}
		\centering
		\includegraphics[width=0.87\textwidth]{Fig_03f.jpg}
		\caption{2019}
	\end{subfigure}%
}\\
\vfill \vspace{1mm}
\makebox[0.90\linewidth][r]{%
	\begin{subfigure}[b]{.40\textwidth}
		\centering
			\includegraphics[width=0.87\textwidth]{Fig_03g.jpg}
		\caption{2020}
	\end{subfigure}%
	\begin{subfigure}[b]{.40\textwidth}
		\centering
		\includegraphics[width=0.87\textwidth]{Fig_03h.jpg}
		\caption{2021}
	\end{subfigure}%
	\begin{subfigure}[b]{.40\textwidth}
		\centering
		\includegraphics[width=0.87\textwidth]{Fig_03i.jpg}
		\caption{2022}
	\end{subfigure}%
}\\
\vfill \vspace{1mm}
\makebox[0.90\linewidth][r]{%
	\begin{subfigure}[b]{.40\textwidth}
		\centering
			\includegraphics[width=0.87\textwidth]{Fig_03j.jpg}
		\caption{2023}
	\end{subfigure}%
	\begin{subfigure}[b]{.40\textwidth}
		\centering
		\includegraphics[width=0.92\textwidth]{Fig_03k.jpg}
		\caption{Google Maps (2023)}
	\end{subfigure}%
	\begin{subfigure}[b]{.40\textwidth}
		\centering
		\includegraphics[width=0.92\textwidth]{Fig_03l.jpg}
		\caption{Google Earth (2023)}
	\end{subfigure}%
}\caption{Landsat 8-9 OLI/TIRS images in natural RGB colours: lake changes for 10 years (2014-2023).}\label{fig:03}
\end{figure}

%----------- subsection ----------->
\subsection{Methods}

The GRASS GIS modules, 'i.cluster' and 'i.maxlik' enable to perform a key procedure converting raster images into classified objects \cite{Neteleretal2008}. The approach of unsupervised classification by the GRASS GIS is based on discriminating the pixels on the image based on their spectral signatures. Upon classification, the areas identified as different land cover classes are represented as geometrical primitives, contoured areas that represent the extent and topological structures of the distinct land cover classes \cite{NetelerMitasova2008}. In the same way, classification enables to treat pixels as different  objects according to the spectral signatures for land cover types grouping them into polygonal segments based on the defined parameters using a clustering algorithm. 

%----------- subsection ----------->
\subsection{Subsection}

%%%%%%%%%%%%%%%%%%%%%%%%%%%%%%%%%%%%%%%%%%
\section{Results}

%%%%%%%%%%%%%%%%%%%%%%%%%%%%%%%%%%%%%%%%%%
\section{Discussion}

%%%%%%%%%%%%%%%%%%%%%%%%%%%%%%%%%%%%%%%%%%
\section{Conclusions}

The principle of using GRASS GIS modules ' i.maxlik ' and 'i.cluster' for creating clusters and image classification, the usage of the scripts by ' i.composite' for color composites from the Landsat bands and the example cases of the used workflow are presented in the scripts. The implemented methods of image analysis are presented together with the examples of GRASS GIS usage in image processing showing identified levels of lakes accurately discriminated from the other land cover types and sand areas of the Sahara desert (Grand Erg Oriental). Maintained geometric shape and topology of classes, automated recognition of pixel’ colours are presented by the GRASS GIS. 

%%%%%%%%%%%%%%%%%%%%%%%%%%%%%%%%%%%%%%%%%%
\authorcontributions{Supervision, conceptualization, methodology, software, resources, funding acquisition, and project administration, O.D.; writing---original draft preparation, methodology, software, data curation, visualization, formal analysis, validation, writing---review and editing, and investigation, P.L. All authors have read and agreed to the published version of the manuscript.}

\funding{The publication was funded by the Editorial Office of XXXX, Multidisciplinary Digital Publishing Institute (MDPI), by providing 100\% discount for the APC of this manuscript. This project was supported by the Federal Public Planning Service Science Policy or Belgian Science Policy Office, Federal Science Policy - BELSPO (B2/202/P2/SEISMOSTORM).}

\institutionalreview{Not applicable.}

\informedconsent{Not applicable.}

\dataavailability{Not applicable} 

\acknowledgments{The authors thank the anonymous reviewers of Land for reading, suggestions and comments that improved an earlier version of this manuscript.}

\conflictsofinterest{The authors declare no conflict of interest.} 

\abbreviations{Abbreviations}{
The following abbreviations are used in this manuscript:\\

\noindent 
\begin{tabular}{@{}ll}
MDPI & Multidisciplinary Digital Publishing Institute\\
DOAJ & Directory of open access journals\\
LD & Linear dichroism
\end{tabular}
}

\section[\appendixname~\thesection]{}
All appendix sections must be cited in the main text. In the appendices, Figures, Tables, etc. should be labeled, starting with ``A''---e.g., Figure A1, Figure A2, etc.

%%%%%%%%%%%%%%%%%%%%%%%%%%%%%%%%%%%%%%%%%%
\begin{adjustwidth}{-\extralength}{0cm}
%\printendnotes[custom] % Un-comment to print a list of endnotes

\reftitle{References}
\nocite{*}

%=====================================
% References, variant A: external bibliography
%=====================================
\bibliography{Algeria.bib}

%%%%%%%%%%%%%%%%%%%%%%%%%%%%%%%%%%%%%%%%%%
\end{adjustwidth}
\end{document}
